\begin{proposition}[P4.19]\hypertarget{tag-02WQ}{}\label{night:P4.19}\label{stmt:P4.19}
\StatementRecord{night:P4.19}{proved}{02WQ}{\SeriesRef{night:P4.16}, \SeriesRef{night:P4.18}}

The germ $y(x^2-zy^2)=0$ has ramified conductor, of multiplicity
two on each normalization component, but its affine nearby
fibre is $\mathbb C^*\times\mathbb C$. Ramification alone
therefore does not give the spherical intersection obstruction
of P4.16--P4.18. This is a local statement, not a compact
counterexample to 4.2.
\end{proposition}
\begin{proof}

\begin{enumerate}[label={$\langle1\rangle$\arabic*.},leftmargin=*]
\item One normalization component is the plane $y=0$; the other is
the Whitney normalization $(u,v)\mapsto(uv,v,u^2)$.
The singular locus is the smooth $z$-axis. The plane conductor
is $x^2=0$, obtained by restricting the other component's
equation. On the Whitney component the intrinsic conductor
is $v\mathbb C\{u,v\}$; multiplication by the added plane
equation $y=v$ gives $v^2\mathbb C\{u,v\}$. Directly, the
union ring is the fibre product of the component rings over
$\mathbb C\{x,z\}/(x^2)$, giving exactly these conductor ideals.
The reduced conductor cover is the identity sheet from the
plane together with $u\mapsto u^2$ from the Whitney component.
The latter is ramified.

\item On $y(x^2-zy^2)=1$ we have $y\ne0$, and the inverse maps
\[
 (x,y,z)\longmapsto(y,x),\qquad
 (y,x)\longmapsto\left(x,y,\frac{yx^2-1}{y^3}\right)
\]
identify this fibre with $\mathbb C^*\times\mathbb C$.
Its $\pi_2$ and $H_2$ vanish. Hence every compact spherical
class in an open subset has zero square after inclusion,
and therefore has zero square already in that open subset.

\end{enumerate}
\end{proof}